% !TEX program = lualatex
\documentclass[11pt,a4paper,sans]{moderncv}

% Настройки для LuaLaTeX
\usepackage{fontspec}
\defaultfontfeatures{Renderer=Basic, Ligatures=TeX}
\newfontfamily\cyrillicfonttt{CMU Typewriter Text}
\newfontfamily\cyrillicfont{CMU Sans Serif}
\setmainfont{CMU Sans Serif}
\setsansfont{CMU Sans Serif}
\setmonofont{CMU Typewriter Text}

\usepackage{polyglossia}
\setdefaultlanguage{russian}
\setotherlanguages{english}

% Настройки ModernCV
\moderncvstyle{casual}
\moderncvcolor{blue}
\nopagenumbers{}

% Настройка полей
\usepackage[scale=0.75]{geometry}

\firstname{Степанова}
\familyname{Анна}
\title{Разработчик Python}
\address{Санкт-Петербург}{}{Россия}
\mobile{+7 (981) 122-43-93}
\email{st.ann0@yandex.ru}
\social[telegram]{slouchdog}
\social[github]{Stepanova-Anna} 

% Начало документа
\begin{document}
\makecvtitle

\section{О себе}
Студентка 3 курса по направлению «Информатика и вычислительная техника», 
активно развиваюсь в направлении Python-разработки. Ищу позицию стажера, 
чтобы применить академические знания и практические навыки (Python, Git, SQL, ООП) 
в реальных проектах. Имею опыт разработки Telegram-бота в рамках производственной 
практики. Отличаюсь аналитическим складом ума, быстрой обучаемостью и 
нацеленностью на результат. Готова решать сложные задачи и расти как 
специалист в сильной команде.

\section{Образование}
\cventry{2023--2027}{Специализация: 09.03.01 «Информатика и вычислительная техника»}{РГПУ им. А. И. Герцена}{Санкт-Петербург}{}{
\begin{itemize}
  \item Дополнительная специализация: «Технологии разработки программного обеспечения и обработки больших данных»
  \item Форма обучения: Очная
\end{itemize}}

\section{Опыт работы}
\cventry{Февраль 2024}{Практикант}{ООО «Галактика»}{}{}{
\begin{itemize}
  \item Разработала Telegram-бота на Python (оценка «отлично»)
\end{itemize}}

\cventry{Сентябрь 2024}{Практикант}{ООО «Галактика»}{}{}{
\begin{itemize}
  \item Участвовала в создании веб-сайта для коммерческого продукта (оценка «отлично»)
\end{itemize}}

\section{Ключевые навыки}
\cvitem{Языки программирования}{Python, Java, C, SQL}
\cvitem{Базы данных}{MySQL}
\cvitem{Инструменты}{Git, ООП, алгоритмы и структуры данных}
\cvitem{Веб-разработка}{Основы разработки веб-приложений, клиент-серверное взаимодействие}
\cvitem{Аналитические способности}{Аналитический склад ума, внимательность к деталям, развитое логическое и стратегическое мышление, опыт обработки и анализа больших данных}
\cvitem{Обучаемость}{Целеустремленность, быстрая обучаемость, усидчивость, готовность к обучению и решению сложных задач, изучение новых технологий и фреймворков}
\cvitem{Командная работа}{Умение работать в команде, подтвержденное опытом участия в хакатонах и IT-соревнованиях}

\section{Сертификаты и достижения}
\cvitem{2024}{Интерактивный тренажер по SQL (Stepik, сертификат с отличием)}
\cvitem{2024}{Безопасность в интернете (Stepik, сертификат с отличием)}
\cvitem{}{1 разряд по шахматам}
\cvitem{}{Участие в хакатонах и IT-соревнованиях}

\end{document}
